\section{Preliminaries}
\label{sec:background}

We recall some basics on Bayesian networks and on quantum computation, respectively referring to~\cite{DarwicheBook} (or~\cite{DarwicheHandbook} for a compact introduction) and to~\cite{Nielsen_book,Watrous_book} for further reading.
We then present quantum Bayesian networks from~\cite{GBN2014}.

\subsection{Classical data and Bayesian Networks}
\label{sec:BN}

Bayesian networks (BNs) are probabilistic graphical models.
Bayesian models provide a formalism for reasoning under conditions of uncertainty or partial knowledge: given a system under study, how \emph{likely} is it that a particular feature is in a particular state?
Every feature of the system is represented by a \emph{random variable}.
For the purpose of modeling, each random variable can be seen as a \emph{name} for an atomic proposition (\eg ``Rain'') which assumes values from a set of states (\eg $\{\true,\false\}$).
The full system is modeled as \emph{a joint probability distribution} over the variables of interest; each element in the sample space represents a possible state.

\subparagraph{Random Variables (\rvs).}

Assuming a countable set of names $X,Y,\dots$, a \emph{random variable} (\rv) is the pair of a name $X$ and a finite set of values $\Val X$.
\emph{W.l.o.g.}, we will assume all random variables to be binary, with $\Val{X}=\{x^\true,x^\false\}$ for a name $X$.
We then silently identify a name $X$ with the \rv $(X, \Val X)$.

We adopt the standard convention of capital letters (\eg $X,Y$) denoting random variables, and lowercase letters (\eg $x,y$) for particular \emph{values} of those variables; $\Pr(x)$ stands for $\Pr(X=x)$.
A finite set of names $\bX=\{X_1, \dots, X_n\}$ defines a ``compound'' \rv where $\Val{\bX}$ is the Cartesian product $\Val{X_1} \times \dots \times \Val{X_n}$; we denote by $\bx$ a tuple in $\Val{\bX}$.

% \subparagraph{Queries to a probabilistic model.}
%
% Given a probabilistic model $\Pr(\bX)$, and some variables of interest $\bY \subset\bX$, a typical query is the marginal probability $\Pr(\bY)$ or $\Pr(\bY \mid \bX=\bz)$, the probability of $\bY$ given evidence $\bz$ for $\bZ$.
% The former is obtained by \emph{summing out} the \rvs which are not relevant; the latter via \textbf{Bayes's rule}, which leads us to compute two \textbf{marginals}, $\Pr(\bY, \bZ)$ and ${\Pr(\bZ)}$, because $\Pr(\bY|\bZ)= \frac{\Pr(\bY, \bZ)}{\Pr(\bZ)}$.

\subparagraph{Bayesian Networks.}

A Bayesian network over a set of \rvs $\bX$ is a DAG whose nodes are the \rvs\ and whose directed edges represent conditional dependencies.
It represents in a compact and factorized way a joint probability distribution $\Pr(\bX)$, by exploiting conditional (in)dependencies in the distribution.
This is achieved by associating with each node a \emph{function}---a conditional probability table---which quantifies the dependencies of each node from its parents.

Formally, a \textdef{Bayesian network} $\bn$ over the set of \rvs $\bX$ is a pair $(\G,\Phi)$ where:
\begin{itemize}
\item $\G$ is a directed acyclic graph (DAG), whose set of nodes is $\bX$;
\item $\Phi$ assigns to each $X\in\bX$ a \emph{conditional probability table (\cpt)} $\phi^X$, for $X$ given its parents.
\end{itemize}

\begin{theorem}[Semantics of BNs~\cite{Pearl86}]
\label{lem:BNsem}
A BN $\bn$ over the set of \rvs $\bX$ defines a unique probability distribution $\Pr(\bX)$ given by the \emph{product} of the CPTs associated to the nodes:
\begin{equation}
\label{eq:BNsem}
\tag{$\ast$}
\Pr(\bX) = \prod_{X\in \bX} \phi^X
\end{equation}
\end{theorem}

\begin{remark}[Markov condition]
Given a Bayesian network over $\bX$, each variable $X$ is assumed to be independent of its non-descendants, given its parents.
A probability distribution $\Pr(\bX)$ can be expressed as (\ref{eq:BNsem}) if and only if it satisfies such a Markov condition.
\end{remark}

The \emph{marginal distribution} over the \emph{variables of interest} $\bY \subseteq \bX$ is obtained by summing-out the irrelevant variables:
$\Pr(\bY)= \sum_{\bX\difset\bY} \Pr(\bX)$.
 
\subparagraph{From CPTs to Factors.}

The strength of BNs is to represent large probability distributions compactly, and to allow for inference algorithms that do not construct the full joint distribution explicitly.
CPTs alone are insufficient for this purpose, motivating the definition of a class of functions called \emph{factors}, a general mathematical abstraction representing both the initial parameters (the CPTs) and the partial results generated during the inference.
Formally, a \emph{factor} $\someqfact$ over the \rvs\ $\bX$ is a function mapping each $\bx\in\setofrv{\bX}$ to a non-negative real:
\begin{equation}
\label{eq:factor}
\someqfact:\setofrv{\bX}\to\nnReal
\end{equation}
A CPT $\theta$ for $X$ given $\bY$ is a factor over $X,\bY$ which satisfies $\sum_x \theta({x,\by}) = 1$ for all $\by\in\setofrv{\bY}$.

Unlike CPTs, factors are closed under both \emph{product and sum-out} operations.
This is crucial for algorithms (\eg\ Variable Elimination) that rely on pushing the sum (marginalization) to smaller components to reduce the cost of computations.
The so generated partial results are factors, but not CPTs.
Moving to factors is also necessary to incorporate evidence.

\subparagraph{Sharing variables.}

We refer to~\cite{DarwicheHandbook} for the formal definitions of sum and product of factors, which are the natural extensions of the corresponding operations for probabilities.
Tractability relies on the properties of these operations, and specifically on two key aspects: the product of factors \emph{inherently shares variables}, and product and sum \emph{distribute} under suitable conditions. Distributivity allows for pushing the sum on smaller components, reducing the size of computations.
Let us briefly discuss how sharing is built in the definition of product, and what this means. 
Assume we want to compute the distribution $\Pr(X,Y,Z)$ defined as the product $ \Pr(Z) \Pr(X \mid Z) \Pr(Y \mid Z)$.
This amounts to computing $2^3$ entries:
\begin{equation*}
\Pr(x,y,z) = \Pr(z) \Pr(x \mid z) \Pr(y \mid z) \quad \text{ for } x\in \setofrv{X}, y\in \setofrv{Y}, z\in \setofrv{Z}
\end{equation*}
Observe how the values corresponding to the same \rv\ $Z$ are shared: we compute $\Pr(z^\true) \Pr(x \mid z^\true) \Pr(y \mid z^\true)$ and $\Pr(z^\false) \Pr(x \mid z^\false) \Pr(y \mid z^\false)$, but never $\Pr(z^\true) \Pr(x \mid z^\true) \Pr(y \mid z^\false)$!

It is important to stress that the definition of the factors product is starkly different with that of the tensor product of matrices, which is the natural product for quantum systems.
We compare the two operations in \cref{app:section2}---the \emph{computational cost} being also starkly different (as discussed in~\cite[Example~9.7]{popl24}).
% In this difference lies the efficiency of a factors-based semantics w.r.t. to a categorical [Jacobs and Zanasi 2020] or relational [Ehrhard et al. 2014; Ehrhard and Tasson 2019] one, where a central role is played by a product ⊗ which behaves as the tensor product of matrices.

\subsection{Quantum data}

An (isolated) quantum system is associated with a complex vector space equipped with an inner product $\braket{\cdot | \cdot}$, \ie\ a \emph{Hilbert space} $\HH$ called the \definitive{state space} of the system.
The simplest quantum system is the \emph{qubit}, which has a two-dimensional state space corresponding to $\C^2$.

The \textbf{states} of a quantum system can be described in two mathematically equivalent ways, either as \emph{vectors} \emph{in} the space $\HH$ or as \emph{linear operators} acting \emph{on} the space $\HH$.
The latter formulation---standard in physics---shines in the description of quantum systems whose state is not known, as they allow to present a state as a probability distribution of states.

\begin{example}
\label{ex:preparation}
Intuitively, a state corresponds to a \emph{preparation}, \ie\ an action performed to set up an experiment.
Given two states $\rho_0$ and $\rho_1$, there must be a state $r\rho_0 + (1-r)\rho_1$ for any $r\in[0,1]$.
Indeed, a valid way to build a state from two preparations $\rho_0$ and $\rho_1$ is the preparation ``flip a fair coin: if heads, continue as the preparation $\rho_0$, otherwise continue as the preparation $\rho_1$'', yielding $\rho = \tfrac12\rho_0 + \tfrac12\rho_1$.
A probability distribution of states (such as $\rho$) also arises naturally to express \emph{partial knowledge} (or \emph{degrees of belief}) about the state of a system.
For example, $\rho$ above expresses we do not know if the state of the system is $\rho_0$ or $\rho_1$.
\end{example}

Before giving further details, we need to recall some notions from linear algebra.

\subparagraph{Hilbert Spaces.}

A finite $n$-dimensional Hilbert space $\HH$ is a complex vector space (isomorphic to) $\C^n$ with an inner product $\braket{\cdot |\cdot}$.
Column vectors are written $\ket{\psi}$, with $\psi$ a label, and their conjugate-transpose $\bra{\psi} \defeq \ket{\psi}^\dagger$. 	
\Cref{fig:algebra} (from~\cite{Nielsen_book}) summarizes relevant notations.

\begin{table}
\begin{tabular}{|>{\centering$}m{3cm}<{$}|m{10cm}|}
\hline
\textbf{Notation} & \textbf{Description}
\\
\hline
z^* & Complex conjugate of the complex number $z$.
\\
\hline
\ket{\psi} & Vector. Also known as a \textit{ket}.
\\
\hline
\bra{\psi} & Vector dual to $\ket{\psi}$. Also known as a \textit{bra}.
\\
\hline
\braket{\varphi | \psi} & Inner product between the vectors $\ket{\varphi}$ and $\ket{\psi}$.
\\
\hline
\ket{\varphi} \otimes \ket{\psi} & Tensor product of $\ket{\varphi}$ and $\ket{\psi}$. 
\\
\hline
\ket{\varphi \psi} & Abbreviated notation for $\ket{\varphi} \otimes \ket{\psi}$.
\\
\hline
\M^T & Transpose of the matrix $\M$.
\\
\hline
\M^{\dagger} & Adjoint (\aka\ Hermitian conjugate, conjugate-transpose) of $\M$:

$
\left[\begin{smallmatrix}
a\phantom{^*} & b \\ c\phantom{^*} & d
\end{smallmatrix}\right]^{\dagger}
=
\left[\begin{smallmatrix}
a^* & c^* \\ b^* & d^*
\end{smallmatrix}\right]
$.
\vspace*{.2em}
\\
\hline
\braket{\varphi | \M | \psi} & Inner product between $\ket{\varphi}$ and $\M \ket{\psi}$.

Equivalently, inner product between $\M^{\dagger} \ket{\varphi}$ and $\ket{\psi}$.
\\
\hline
\end{tabular}
\caption{Main linear algebra notations.}
\label{fig:algebra}
\end{table}

Given a Hilbert space $\HH$, $\L(\HH)$ denotes the set of \definitive{linear operators} \emph{on} $\HH$ (\ie\ from $\HH$ to $\HH$), which we 
identify with their matrix representations.
The \definitive{adjoint} $\M^\dagger$ of a matrix $\M$ is its conjugate-transpose.
We denote by $\id[\L(\HH)]$ the identity matrix in $\L(\HH)$.
The \definitive{trace} of a square matrix $\M = (\M_{ij})_{ij}\in \C^{n\times n}$ is $\tr{\M} = \sum_i \M_{ii}$.

A relevant subclass of $\L(\HH)$ is the set of positive operators $\Pos{\HH}$, playing the role of ``non‑negative real numbers''.
A linear operator $\M$ is \textbf{positive} %(or positive semi-definite)
if $\braket{\psi | \M | \psi} \in \nnReal$ for any $\ket \psi \in \HH$.

\begin{remark}
\label{rem:pos}
In the 1-dimensional space $\C^1$, the set of positive matrices in $\L(\C)=\C^{1\times 1}=\C$ is $\Pos{\C}=\nnReal$.
\end{remark}

\subparagraph{Quantum states as vectors in $\HH$.}

A common way to represent the states of a quantum system is by unit vectors in $\HH$, \ie\ vectors $\ket{\psi}$ such that $\braket{\psi | \psi} = 1$.
Considering a qubit, the vectors $\ket{0} = \left[\begin{smallmatrix}1\\0\end{smallmatrix}\right]$ and $\ket{1} = \left[\begin{smallmatrix}0\\1\end{smallmatrix}\right]$ form an orthonormal basis for its state space $\C^2$: an arbitrary state vector can be written $\ket{\psi} = \alpha \ket{0} + \beta \ket{1}$.
Intuitively, the states $\ket{0}$ and $\ket{1}$ of a qubit are analogous to the two values $0$ and $1$ which a \emph{bit} may take.
The way a qubit differs is that \emph{superpositions} (\ie\ linear combinations) of these two states also exist.
Some important states are $\ket{+} = \frac{1}{\sqrt{2}}(\ket{0} + \ket{1})$ and $\ket{-} = \frac{1}{\sqrt{2}}(\ket{0} - \ket{1})$, which form the Hadamard basis.

\subparagraph{Quantum states as density operators acting on $\HH$.}

Another usual way in physics to describe the state of a quantum system is by a \emph{positive operator $\rho$ with trace $1$}---called a \definitive{density operator}---that \emph{acts} on the state space $\HH$ of the system.
We denote by $\D(\HH) \subset \Pos{\HH}$ the set of density operators acting on $\HH$.
A vector $\ket{\psi}$ is represented by the matrix $\ket{\psi}\bra{\psi}$, called a \emph{pure state}.
If a quantum system is in the state $\rho_{i}$ with probability $p_i$, the density operator for the system is $\sum_{i} p_{i} \rho_{i}$.
This formalizes \Cref{ex:preparation}: if we can prepare qubits in pure states, we can also prepare qubits in any probability distribution of states.

\subparagraph{Composite systems.}

The \emph{state space} of a composite physical system is the tensor product of the state spaces of its components.
Given several systems, with the states $(\rho_i)_{i\in I}$, the \emph{state} of the composite system is $\bigotimes_{i\in I}\rho_i$.
Consider the Hilbert spaces $\HH_1$ with $\{\ket{\psi_i}\}_{i\in I}$ one of its bases, and $\HH_2$ with $\{\ket{\phi_j}\}_{j\in J}$ one of its bases.
Their tensor product $\HH_1 \otimes \HH_2$ is the Hilbert space with basis $\{\ket{\psi_i} \otimes \ket{\phi_j}\}_{(i,j)\in I\times J}$.
A short notation for $\ket{\psi}\otimes\ket{\phi}$ is $\ket{\psi\phi}$.
A pure quantum state in $\HH_1 \otimes \HH_2$ is \emph{separable} when it can be expressed as the tensor product of two vectors of $\HH_1$ and $\HH_2$.
Otherwise, it is \emph{entangled}, as for example the Bell state $\frac{1}{\sqrt{2}}(\ket{00} + \ket{11})$.

The \definitive{partial trace} $\ptr{\HH_1}$ of an operator in $\L(\HH_1\otimes\HH_2)$ is an operator in $\L(\HH_2)$ defined as $\ptr{\HH_1}\left(\M\otimes\M'\right)=\tr{\M}\M'$ for $\M\in\L(\HH_1)$ and $\M'\in\L(\HH_2)$, and extended to the general case by linearity.
The partial trace allows to ``\emph{forget}'' or ``\emph{discard}'' part of a composite system.

\subparagraph{Quantum operations (\aka\ quantum channels).}

States $\M\in\D(\HH)$ transform as $\M' = \E(\M)$, where the map $\E: {\D}(\mathcal{H}) \to {\D}(\mathcal{H'})$, called \definitive{quantum operation} or \definitive{quantum channel}, is:
\begin{enumerate}
\item\label{item:qo_cp}
\emph{Completely Positive}:
$\E$ sends positive operators to positive operators, and do so also when acting on a sub-part of a larger, potentially \emph{entangled}, composite system; \ie\ any ``extension'' $\idop\otimes \Eop$ (with $\idop$ the identity) \emph{carries a positive operator to a positive one}.
\item\label{item:qo_tp}
\emph{Trace-Preserving}:
if $\tr{\rho}=1$ then $\tr{\E(\rho)}=1$.
%\item
%\emph{Trace Non-increasing:}
%The trace of the resulting state $\tr{\E(\rho)}$ must be between 0 and 1, representing the \textbf{probability} that the physical process occurred.
\item\label{item:qo_l}
\emph{Convex-Linear}:
$\E$ acts on a mixture of states by acting on each component individually, \ie\ $\Eop\left(\sum_{i\in I}p_{i}\M_{i}\right) = \sum_{i\in I}p_{i}\Eop(\rho_{i})$.
\end{enumerate}
By \ref{item:qo_cp} and \ref{item:qo_tp}, $\Eop$ (or more generally $\idop\otimes\Eop$) sends a density operator to a density operator.

\begin{example}
Two fundamental ways to modify a quantum state are \emph{unitary transformations}, describing the evolution of an isolated system, and \emph{measurements}.
A matrix $U\in\L(\HH)$ is unitary when $UU^\dagger = U^\dagger U = \id$.
The corresponding quantum channel is $\E:\rho \mapsto U \rho U^\dagger$.
Similarly, measurement operators $M_m$ act on a state $\rho$ as $M_m \rho M_m^\dagger$.
One more example of quantum operation is the partial trace.
\end{example}

\subparagraph{Quantum Instruments.}

Quantum channels describe the evolution of quantum states, \emph{ignoring classical information}---such as measurement outcomes---which plays a central role in ``real'' (or thought) experiments, as in \Cref{fig:Bell}.
A generalization of channels are \emph{quantum instruments}, mathematical devices that model ``real'' devices by specifying a post-measurement state per classical outcome.
This provides a framework where measurements are operations producing both quantum state updates and classical data.
A \textbf{quantum instrument} is an indexed collection $\{\E_a\}_a$ of completely positive linear maps that sum to a trace preserving map (\ie\ to a quantum channel).
Intuitively, the labels $a$ are the \emph{classical} outcomes ($a^\true/a^\false$) of a measurement---think $a\in \setofrv{A}$.
An instrument $\{\mathcal{E}_{a}\}_{a}$ induces a quantum channel $\mathcal{E}(\rho) =\sum_{a} \mathcal{E}_{a}(\rho)$, with $\Pr(a \mid \rho) = \tr{\mathcal{\E}_{a}(\rho)}$ providing the outcome probabilities.
We refer \eg\ to~\cite[page~112]{Watrous_book} for details.

\begin{remark}[Pure quantum and pure classical]
\label{rem:qi_pr}
The two limit cases yield standard notions.
\begin{itemize}
\item
Quantum instruments with no classical outcome (\ie\ with a single label $a$) are simply quantum operations.
\item
Quantum instruments from $\C$ to $\C$ are (isomorphic to) probability distributions.
% Indeed, a linear map $\E_\bx:\C\to\C$ corresponds to a scalar $p_\bx$, with $p_\bx\in\nnReal$ when $\E_\bx$ is (completely) positive.
% The trace condition on $\sum_\bx \E_\bx$ implies $\sum_\bx p_\bx = 1$.
\end{itemize}
\end{remark}

\subsection{Quantum Bayesian Networks (Instrument-based)}
\label{sec:qbn_instrument}

The generalization of Bayesian networks in~\cite{GBN2014} is based on the notion of quantum instruments. In \Cref{sec:qbn}, we will provide a definition of quantum Bayesian Network in terms of \qfactor s---the original definition (to which ours is equivalent) is not necessary to the reader. 

In this background section, we briefly review the key ideas in~\cite{GBN2014}, and what are the issues with the semantics there. 
The causal structure is described by a DAG (as in \cref{fig:Bell}); to include quantum sources of causality, the set of nodes is extended with new unobserved nodes, corresponding to quantum systems.
To each node of the generalized DAG is associated a \emph{family of quantum instruments} rather than a CPT, as we illustrate in \cref{ex:Alice_instrument}.

\begin{remark}
Recall that to each BN node is associated a CPT, \ie\ a \emph{family of probability distributions} indexed by values.
Intuitively, quantum instruments generalize probability distributions, so a \emph{family of quantum instruments} indexed by values generalizes a CPT.
\end{remark}

\begin{example}[Alice's instruments]
\label{ex:Alice_instrument}
Referring to \Cref{fig:Bell}, we want to model that Alice (the node $A$) performs the following operations, depending on the outcome of $X$.
Given $x^t$, Alice measures the qubit $q_1$ on the $(\ket{0},\ket{1})$ basis, returning $a^t$ if she obtains $\ket{0}$ and $a^f$ if not.
Given $x^f$, Alice measures $q_1$ on the $(\ket{-},\ket{+})$ basis, returning $a^t$ if she obtains $\ket{+}$, $a^f$ if not.

To model this setting, we associate to $A$ a \emph{family} of instruments, indexed by $\Val{X}$: $\{\{\Eop_{a^t,x^t},\Eop_{a^f,x^t}\},\{\Eop_{a^f,x^f},\Eop_{a^f,x^f}\}\}$.
The instrument given $x^t$ is $\{\Eop_{a^t,x^t},\Eop_{a^f,x^t}\}$ (from $\L(\HH_{q_1})$ to $\C$) with $\Eop_{a^t,x^t} : \M \mapsto \ket{0}\bra{0}\M\ket{0}\bra{0}$ and $\Eop_{a^f,x^t}: \M \mapsto \ket{1}\bra{1}\M\ket{1}\bra{1}$.
The instrument given $x^f$ is $\{\Eop_{a^f,x^f},\Eop_{a^f,x^f}\}$, % from $\L(\HH_{q_1})$ to $\C$, 
with $\Eop_{a^t,x^f} : \M \mapsto \ket{+}\bra{+}\M\ket{+}\bra{+}$ and $\Eop_{a^f,x^f}: \M \mapsto \ket{-}\bra{-}\M\ket{-}\bra{-}$.
\end{example}

%\begin{remark}[No-Cloning and edges interpretation]
%A defining feature of quantum computing is the \emph{No-Cloning Theorem}, implying that qubits cannot be duplicated.
%As a result, qubits cannot be stored in multiple locations or broadcast to multiple receivers.
%\end{remark}

\subparagraph{Semantics of \QBN s.} 

The interpretation of a \QBN\ $\somebayes$ over a set of classical variables $\bX$ and a set of quantum systems $\bQ$ is defined in~\cite{GBN2014} as a suitable product of the instruments associated to the nodes.
Similarly to BNs, the semantics $\sem{\somebayes}$ of $\somebayes$ is a probability distribution over $\bX$, reflecting the fact that quantum states cannot be observed directly---we can only observe the classical outcomes of measurements.
This is exactly what happens in the Bell set-up (\Cref{ex:bell}), where we compute $\Pr(a,b\mid x,y)$. 

\subparagraph{Lack of compositionality.}
\label{sec:lack}

While intuitive and natural, the approach via quantum instruments has a limit in \emph{the lack of compositionality}.
Indeed, the interpretation of a \QBN\ $\somebayes$ in~\cite{GBN2014} needs to follow a specific order (see ~\cite[page~12]{GBN2014}).
The definition of an instrument does not adapt well to interpreting arbitrary sub-components (as in \cref{fig:compositional}).
Given a partition of $\somebayes$ into $\somebayes_1$ and $\somebayes_2$, even if we can define $\sem{\somebayes_1}$ and $\sem{\somebayes_2}$, we do not have a natural composing function yielding $\sem{\somebayes_1} \bullet \sem{\somebayes_2}=\sem{\somebayes}$.
There are two distinct issues.
\begin{enumerate}
\item
Consider \Cref{fig:compositional} with its \textcolor{red}{left} and \textcolor{blue}{right} sub-graphs.
The \textcolor{red}{left} sub-graph should be interpreted by an instrument $\{\Eop^l_x \mid x \in\setofrv{X}\}$ with $\Eop^l_x : \HH_2\to \HH_1$, and the \textcolor{blue}{right} sub-graph by $\Eop^r : \HH_1\to \HH_2$.
The interpretation of the full graph is an instrument $\{\Eop_x \mid x \in\setofrv{X}\}$ with $\Eop_x:\C\to\C$. 
However, the operations on instruments do not immediately yield from $\Eop^l_x$ and $\Eop^r$ a function of type $\C\to\C$.
% However, it is impossible to get a function of type $\C\to\C$ using the $\Eop^l_x$, $\Eop^r$ and only identity functions, tensor products and the usual composition of functions (which are the operations used in \cref{lem:qbn_qi_pb}).
\item
A further issue is the impossibility (in general) of dealing by components with hidden variables, exactly as happens for CPTs, which are closed by product but not by sum.
\end{enumerate}

\begin{remark}[Marginals and hidden variables]
\label{rem:marginals}
A technical but crucial point when considering compositionality is the variables of interest---or dually, the hidden variables.
In \cref{fig:compositional}, the shaded node $Y$ corresponds to a hidden (or unobserved) variable.
What we want here is the marginal $\Pr(x)$, which can be obtained from $\Pr(x,y)$ by summing-out the irrelevant variable: $\sum_{y \in \setofrv{Y}} \Pr(x,y)$.
When computing the semantics of the \textcolor{blue}{right} component, it would be desirable to already sum-out $Y$.
This kind of operation is regularly performed in Bayesian inference algorithms to reduce the size of computations.
\end{remark}

Addressing both issues with compositionality motivates the framework of \emph{quantum factors} that we introduce in the next section.
In the classical setting, abstracting CPTs into \emph{factors} gives a status to partial calculations. We will follow a similar path.
